\documentclass[11pt]{amsart}
\usepackage[a4paper,margin=1.05in]{geometry}
\usepackage[T1]{fontenc}
\usepackage{lmodern,microtype,amsmath,amssymb,amsthm,mathtools}
\usepackage[hidelinks]{hyperref}
\hypersetup{pdftitle={The maximal Kemeny constant of an n-vertex graph},pdfauthor={Yitzchak Shmalo}}
\setlength{\emergencystretch}{1.5em}
\numberwithin{equation}{section}
\newtheorem{theorem}{Theorem}[section]
\newtheorem{proposition}[theorem]{Proposition}
\newtheorem{lemma}[theorem]{Lemma}
\newtheorem{corollary}[theorem]{Corollary}
\theoremstyle{remark}
\newtheorem{remark}[theorem]{Remark}
\newcommand{\E}{\mathbb E}
\newcommand{\trel}{\tau_{\mathrm{rel}}}
\newcommand{\dist}{\operatorname{dist}}
\newcommand{\diam}{\operatorname{diam}}
\title[The maximal Kemeny constant]{The maximal Kemeny constant of an $n$-vertex graph}
\author{Yitzchak Shmalo}
\address{Einstein Institute of Mathematics, The Hebrew University of Jerusalem, Jerusalem, Israel}
\email{yitzchak.shmalo@gmail.com}
\date{10 October 2026}
\subjclass[2020]{Primary 60J10, 05C81; Secondary 05C35, 05C50}
\keywords{Kemeny's constant, random walk, effective resistance, extremal graph, barbell graph, Tur\'an graph}
\begin{document}
\begin{abstract}
For simple random walk on a connected simple graph with $n$ vertices, Kemeny's constant is at most $n^3/54+\frac{37}{8}n^{5/2}$. Balanced barbells show that the leading term is sharp, so the maximum of Kemeny's constant over such graphs is asymptotic to $n^3/54$. This proves the Kemeny-constant assertion of Aldous--Fill Open Problem~6.13. Since the relaxation time never exceeds Kemeny's constant, the same bound gives the relaxation-time assertion of that problem, which was proved by Aksoy, Chung, Tait and Tobin. The proof is elementary. It combines the effective-resistance formula for Kemeny's constant with a contraction of the neighbourhoods of high-degree vertices and a two-mass inequality on the edges of a spanning tree. The same argument shows that among $K_{r+1}$-free graphs, and among graphs of chromatic number at most $r$, both maxima equal $(1-1/r)n^3/54+O(n^{5/2})$. In particular the triangle-free and the bipartite constants are $1/108$.
\end{abstract}
\maketitle

\section{Introduction}\label{sec:intro}
Let $G$ be a connected simple graph with vertex set $V$, $|V|=n\ge2$, degrees $d_v$, and $W=\sum_vd_v=2|E(G)|$. Simple random walk $(X_t)$ on $G$ has transition probabilities $P(u,v)=1/d_u$ when $uv\in E(G)$ and $P(u,v)=0$ otherwise, and stationary distribution $\pi_v=d_v/W$. The eigenvalues of $P$ are real, $1=\lambda_1>\lambda_2\ge\cdots\ge\lambda_n\ge-1$, and we write
\begin{equation}\label{eq:Kdef}
K(G)=\sum_{j=2}^n\frac1{1-\lambda_j},\qquad \trel(G)=\frac1{1-\lambda_2}.
\end{equation}
With $T_v=\min\{t\ge0:X_t=v\}$, Kemeny's constant $K(G)$ equals $\sum_v\pi_v\E_uT_v$ for every starting vertex $u$ (Lemma~\ref{lem:formulas}). In the notation of Aldous and Fill, $K(G)=\tau_0$ and $\trel(G)=\tau_2$ for the walk run in continuous time with unit jump rate, whose mean hitting times are those of the discrete-time walk. Every term of the sum in \eqref{eq:Kdef} is positive, so $\trel(G)\le K(G)$.

Aldous and Fill \cite[Chapter~6, Open Problem~6.13]{AF} ask for proofs of
\begin{equation}\label{eq:AF}
\max_GK(G)\sim\frac{n^3}{54},\qquad \max_G\trel(G)\sim\frac{n^3}{54},
\end{equation}
the maxima being over connected simple graphs on $n$ vertices. Aksoy, Chung, Tait and Tobin \cite{ACTT} proved the second assertion. The first is proved here, with an explicit error term.

\begin{theorem}\label{thm:main}
Every connected simple graph $G$ on $n\ge2$ vertices satisfies
\begin{equation}\label{eq:upper}
\trel(G)\le K(G)\le\frac{n^3}{54}+\frac{37}{8}\,n^{5/2}.
\end{equation}
For every $n\ge6$ there is a connected simple graph $B_n$ on $n$ vertices, formed by two complete graphs of order $\lfloor n/3\rfloor$ joined by a path, with
\begin{equation}\label{eq:lower}
K(B_n)\ge\trel(B_n)\ge\frac{n^3}{54}-\frac{n^2}{18}-\frac{5n}{54}.
\end{equation}
Consequently $\max_GK(G)\sim\max_G\trel(G)\sim n^3/54$.
\end{theorem}

The upper bound uses only how many edges a set of vertices can span, and so it improves under restrictions of that kind.

\begin{theorem}\label{thm:clique}
Fix an integer $r\ge2$. Every connected $K_{r+1}$-free graph $G$ on $n\ge2$ vertices satisfies
\begin{equation}\label{eq:clique-upper}
\trel(G)\le K(G)\le\Bigl(1-\frac1r\Bigr)\frac{n^3}{54}+\frac{37}{8}\,n^{5/2}.
\end{equation}
For every $n\ge6r$ there is a connected graph of chromatic number at most $r$ on $n$ vertices with
\begin{equation}\label{eq:clique-lower}
K(G)\ge\trel(G)\ge\Bigl(1-\frac1r\Bigr)\frac{n^3}{54}-\frac{(r-1)n^2}{9}.
\end{equation}
Hence, over connected $K_{r+1}$-free graphs and over connected graphs of chromatic number at most $r$, the maxima of $K$ and of $\trel$ are both $(1-1/r)n^3/54+O(n^{5/2})$. In particular, over triangle-free graphs and over bipartite graphs both maxima are $n^3/108+O(n^{5/2})$.
\end{theorem}

The lower bounds come from barbells, two copies of a dense graph joined by a long path. Exact formulas for the Kemeny constant of barbells are given in \cite{Breen,Faught}. The content of Theorem~\ref{thm:main} is the upper bound over all graphs.

The constant $1/54$ has a simple source. By the resistance formula of Lemma~\ref{lem:formulas}, $K(G)$ is a degree-weighted average of effective resistances. In a barbell whose two blocks have $a$ vertices each and whose path has length $\ell$, two masses of degree sum about $a^2$ sit at resistance about $\ell$ from each other, so $K$ is about $a^2\cdot a^2\cdot\ell/(2a^2)=a^2\ell/2$. Under $2a+\ell\approx n$ this is largest at $a=\ell=n/3$, where it equals $n^3/54$. The proof turns this heuristic into an inequality. The vertices of degree at least $\sqrt n$, together with their neighbours, form clusters of small total diameter (Lemma~\ref{lem:clusters}). After the clusters are contracted, the graph is compared with a spanning tree, and each tree edge separates two masses whose contribution is at most $s^2/8$, where $s$ is the number of clustered vertices (Lemma~\ref{lem:twomass}). The tree has about $n-s$ edges, and $(n-s)s^2/8\le n^3/54$.

\section{Kemeny's constant and effective resistance}\label{sec:resistance}
Let $L=\operatorname{diag}(d)-A$ be the Laplacian of $G$, where $A$ is the adjacency matrix, and let $e_u$ denote the coordinate vectors. For distinct $u,v\in V$ the equation $L\varphi=e_u-e_v$ has a solution, unique up to an additive constant, and the effective resistance is $R_{uv}=\varphi_u-\varphi_v$. Put $R_{uu}=0$.

\begin{lemma}\label{lem:formulas}
Let $G$ be a connected simple graph and $u,v\in V$. Then
\begin{enumerate}
\item[(a)] $K(G)=\sum_{x}\pi_x\E_uT_x$;
\item[(b)] $\E_uT_v+\E_vT_u=WR_{uv}$;
\item[(c)] $K(G)=\frac1W\sum_{\{x,y\}}d_xd_yR_{xy}$, the sum running over unordered pairs of distinct vertices;
\item[(d)] $R_{uv}\le\dist_G(u,v)$.
\end{enumerate}
\end{lemma}

These facts are standard \cite[Chapter~3]{AF}; we include short proofs.

\begin{proof}
The matrix $P$ is self-adjoint on $\ell^2(\pi)$, and
\begin{equation}\label{eq:dirichlet}
\sum_x\pi_xf_x\bigl((I-P)f\bigr)_x=\frac1W\sum_{xy\in E(G)}(f_x-f_y)^2 .
\end{equation}
Since $G$ is connected, the eigenvalue $1$ of $P$ is simple and its eigenvectors are the constants. Put $\Pi=\mathbf 1\pi^{\mathsf T}$ and $Z=(I-P+\Pi)^{-1}$. The matrix $I-P+\Pi$ maps $\mathbf 1$ to itself and acts as $I-P$ on the functions with $\pi^{\mathsf T}f=0$, so it has eigenvalues $1$ and $1-\lambda_j$ for $j\ge2$. Thus $Z$ exists, $\operatorname{tr}Z=1+K(G)$, $Z\mathbf 1=\mathbf 1$, and, since $\Pi Z=\Pi$, also $(I-P)Z=I-\Pi$.

(a) Fix $v$ and let $h_x=(Z_{vv}-Z_{xv})/\pi_v$. Then $h_v=0$, and for $x\ne v$ we get $((I-P)h)_x=-((I-P)Z)_{xv}/\pi_v=\pi_v/\pi_v=1$. The function $x\mapsto\E_xT_v$ satisfies the same equations by first-step analysis, and the difference of the two functions vanishes at $v$ and is harmonic for $P$ elsewhere, so it is zero by the maximum principle. Hence $\E_uT_v=(Z_{vv}-Z_{uv})/\pi_v$, and $\sum_v\pi_v\E_uT_v=\operatorname{tr}Z-\sum_vZ_{uv}=\operatorname{tr}Z-1=K(G)$.

(b) Let $g^v_x=\E_xT_v$. For $x\ne v$ we have $(Lg^v)_x=d_x((I-P)g^v)_x=d_x$. The column sums of $L$ vanish, so $(Lg^v)_v=-\sum_{x\ne v}d_x=d_v-W$. Thus $Lg^v=d-We_v$, and $L(g^v-g^u)=W(e_u-e_v)$. Taking $\varphi=(g^v-g^u)/W$ gives $R_{uv}=\varphi_u-\varphi_v=(\E_uT_v+\E_vT_u)/W$.

(c) Averaging (a) over $u$ with weights $\pi_u$ and using (b),
\[
K(G)=\sum_{x,y}\pi_x\pi_y\E_xT_y=\frac12\sum_{x,y}\pi_x\pi_y\bigl(\E_xT_y+\E_yT_x\bigr)=\frac W2\sum_{x,y}\pi_x\pi_yR_{xy}=\frac1W\sum_{\{x,y\}}d_xd_yR_{xy}.
\]

(d) Orient each edge and let $J$ be the vertex-edge incidence matrix, with entry $1$ at the tail and $-1$ at the head of each edge, so that $L=JJ^{\mathsf T}$. If $L\varphi=e_u-e_v$, the edge vector $i=J^{\mathsf T}\varphi$ satisfies $Ji=e_u-e_v$ and $\|i\|^2=\varphi^{\mathsf T}L\varphi=R_{uv}$. If $j$ is any other edge vector with $Jj=e_u-e_v$, then $\langle i,j-i\rangle=\langle\varphi,J(j-i)\rangle=0$, so $\|j\|^2=\|i\|^2+\|j-i\|^2\ge R_{uv}$. The unit flow along a shortest path from $u$ to $v$ has $\|j\|^2=\dist_G(u,v)$.
\end{proof}

\section{The upper bound}\label{sec:upper}
We prove the upper bounds of Theorems~\ref{thm:main} and~\ref{thm:clique} in the following form.

\begin{theorem}\label{thm:capacity}
Let $0<\gamma\le1$ and let $G$ be a connected simple graph on $n\ge2$ vertices such that
\begin{equation}\label{eq:capacity}
2|E(G[U])|\le\gamma|U|^2\qquad\text{for every }U\subseteq V.
\end{equation}
Then
\[
K(G)\le\frac{\gamma n^3}{54}+\frac{37}{8}\,n^{5/2}.
\]
\end{theorem}

Every simple graph satisfies \eqref{eq:capacity} with $\gamma=1$. By Tur\'an's theorem (see \cite{MS}), a $K_{r+1}$-free graph on $m$ vertices has at most $(1-1/r)m^2/2$ edges, and induced subgraphs of a $K_{r+1}$-free graph are $K_{r+1}$-free; hence every $K_{r+1}$-free graph satisfies \eqref{eq:capacity} with $\gamma=1-1/r$.

For a vertex $h$ write $N[h]$ for its closed neighbourhood, $h$ together with its neighbours.

\begin{lemma}\label{lem:clusters}
Let $G$ be a finite simple graph and $D>0$. Let $H=\{v:d_v\ge D\}$, $S=\bigcup_{h\in H}N[h]$, and let $C_1,\dots,C_k$ be the vertex sets of the connected components of $G[S]$. Put $s_i=|C_i|$, $s=|S|=\sum_is_i$ and $\delta_i=\diam G[C_i]$. Then
\[
k\le\frac sD\qquad\text{and}\qquad\sum_{i=1}^k\delta_i\le\frac{5s}D.
\]
\end{lemma}

\begin{proof}
Let $v\in C_i$, and choose $h\in H$ with $v\in N[h]$. The set $N[h]$ lies in $S$ and induces a connected subgraph, and it contains $v$, so $N[h]\subseteq C_i$. In particular every $C_i$ contains a set $N[h]$ with $h\in H$, so $s_i\ge d_h+1>D$, and therefore $k<s/D$ when $k\ge1$.

Fix $i$ and a shortest path $v_0,v_1,\dots,v_{\delta_i}$ in $G[C_i]$. For each index $j\in\{0,5,10,\dots\}$ with $j\le\delta_i$ choose $h_j\in H$ with $v_j\in N[h_j]$; as above, $N[h_j]\subseteq C_i$. If $N[h_j]$ and $N[h_l]$ meet for two such indices $j<l$, then $G[C_i]$ contains a walk from $v_j$ through $h_j$, a common vertex of the two neighbourhoods and $h_l$ to $v_l$ of length at most $4$, which is impossible because $\dist_{G[C_i]}(v_j,v_l)=l-j\ge5$. So these $\lfloor\delta_i/5\rfloor+1$ sets are disjoint subsets of $C_i$, each with more than $D$ vertices. Hence $(\delta_i/5)D<(\lfloor\delta_i/5\rfloor+1)D<s_i$, so $\delta_i<5s_i/D$. Summing over $i$ gives the second inequality.
\end{proof}

\begin{lemma}\label{lem:twomass}
Let $\gamma>0$, let $x,y\ge0$ with $x+y=s$, and let $0\le\alpha\le\gamma x^2$ and $0\le\beta\le\gamma y^2$. Then
\[
\alpha\beta\le\frac{\gamma s^2}8\,(\alpha+\beta).
\]
\end{lemma}

\begin{proof}
If $\alpha\beta=0$ there is nothing to prove, so let $\alpha,\beta>0$; then $x,y>0$. The function $(\alpha,\beta)\mapsto\alpha\beta/(\alpha+\beta)$ is nondecreasing in each variable on $(0,\infty)^2$, its partial derivative in $\alpha$ being $\beta^2/(\alpha+\beta)^2$. Hence
\[
\frac{\alpha\beta}{\alpha+\beta}\le\frac{\gamma^2x^2y^2}{\gamma x^2+\gamma y^2}=\frac{\gamma x^2y^2}{x^2+y^2}\le\frac{\gamma(x+y)^2}8,
\]
where the last step is the identity $(x+y)^2(x^2+y^2)-8x^2y^2=(x-y)^2(x^2+4xy+y^2)$.
\end{proof}

\begin{proof}[Proof of Theorem~\ref{thm:capacity}]
Fix $D>0$ and let $H$, $S$, $C_1,\dots,C_k$, $s$ and $\delta_i$ be as in Lemma~\ref{lem:clusters}. Put $\delta=\sum_i\delta_i$ and $t=n-s$.

\emph{A quotient tree.} Partition $V$ into the parts $C_1,\dots,C_k$ and the singletons $\{v\}$ with $v\notin S$, and write $q(u)$ for the part containing $u$. Let $Q$ be the graph on these $k+t$ parts in which two parts are adjacent when an edge of $G$ joins them. Then $Q$ is connected; let $T$ be a spanning tree of $Q$. By Lemma~\ref{lem:clusters}, $T$ has $k+t-1\le t+s/D$ edges, and it has at most $n-1$ edges.

We claim that $\dist_G(u,v)\le\dist_T(q(u),q(v))+\delta$ for all $u,v\in V$. Let $q(u)=y_0,y_1,\dots,y_m=q(v)$ be the path in $T$, and for each $l$ choose an edge of $G$ from a vertex $p_l\in y_{l-1}$ to a vertex $p_l'\in y_l$. Join $u$ to $p_1$ inside $y_0$, each $p_l'$ to $p_{l+1}$ inside $y_l$, and $p_m'$ to $v$ inside $y_m$. Inside a part $C_i$ such a connection has length at most $\delta_i$, and inside a singleton it has length $0$. The tree path visits each part at most once, so the resulting walk has length at most $m+\delta$. If $q(u)=q(v)$ the bound $\dist_G(u,v)\le\delta$ holds directly.

For a part $y$ let $w_y=\sum_{u\in y}d_u$. Lemma~\ref{lem:formulas}(c),(d) and $\sum_{\{u,v\}}d_ud_v\le W^2/2$ give
\[
K(G)\le\frac1W\sum_{\{u,v\}}d_ud_v\bigl(\dist_T(q(u),q(v))+\delta\bigr)\le\frac1W\sum_{\{y,z\}}w_yw_z\dist_T(y,z)+\frac{W\delta}2 .
\]
An edge $e$ of $T$ lies on the tree path between $y$ and $z$ exactly when $e$ separates $y$ from $z$. Writing $W_e^-$ and $W_e^+$ for the sums of $w_y$ over the two components of $T-e$, we obtain
\begin{equation}\label{eq:treesum}
K(G)\le\sum_{e\in E(T)}\frac{W_e^-W_e^+}{W}+\frac{W\delta}2 .
\end{equation}

\emph{One tree edge.} For a part $y=C_i$ put $a_y=2|E(G[C_i])|$, and for a singleton put $a_y=0$. Then $b_y=w_y-a_y$ is the number of edges of $G$ with exactly one end in $y$. Distinct components of $G[S]$ are not joined by edges, so every edge between two different parts has an end outside $S$, and every vertex outside $S$ has degree less than $D$. Hence
\[
B:=\sum_yb_y\le2\sum_{v\notin S}d_v\le2tD\le2nD .
\]
Fix $e\in E(T)$, let $\alpha$ and $\beta$ be the sums of $a_y$ over the two sides of $e$, let $b^-$ and $b^+$ be the corresponding sums of $b_y$, and put $A=\alpha+\beta$, so that $W=A+B$. Then
\[
W_e^-W_e^+=(\alpha+b^-)(\beta+b^+)\le\alpha\beta+AB+B^2=\alpha\beta+BW .
\]
Let $x^-$ and $x^+$ be the numbers of vertices of $S$ on the two sides of $e$, so $x^-+x^+=s$. By \eqref{eq:capacity} and $\sum s_i^2\le(\sum s_i)^2$, taken over the clusters on one side, we have $\alpha\le\gamma(x^-)^2$ and $\beta\le\gamma(x^+)^2$. Lemma~\ref{lem:twomass} gives $\alpha\beta\le\gamma s^2A/8\le\gamma s^2W/8$, and therefore
\[
\frac{W_e^-W_e^+}W\le\frac{\gamma s^2}8+B .
\]

\emph{Summation.} Insert this into \eqref{eq:treesum}, count the terms $\gamma s^2/8$ with the bound $t+s/D$ on the number of tree edges and the terms $B$ with the bound $n-1$, and use $\delta\le5s/D$:
\[
K(G)\le\frac{\gamma(n-s)s^2}8+\frac{\gamma s^3}{8D}+2n^2D+\frac{5Ws}{2D}.
\]
Now $\gamma\le1$, $s\le n$ and $W<n^2$, and $\max_{0\le s\le n}(n-s)s^2=4n^3/27$, attained at $s=2n/3$. Hence
\[
K(G)\le\frac{\gamma n^3}{54}+\frac{21}{8}\cdot\frac{n^3}D+2n^2D .
\]
The choice $D=\sqrt n$ completes the proof.
\end{proof}

\section{Barbells}\label{sec:barbells}
Let $F$ be a connected graph with a distinguished vertex $o$, and let $\ell\ge1$ be an integer. The barbell $B(F,\ell)$ consists of two disjoint copies of $F$ whose copies of $o$ are joined by a path with $\ell$ edges and $\ell-1$ new internal vertices. It has $2|V(F)|+\ell-1$ vertices.

\begin{lemma}\label{lem:barbell}
For every connected graph $F$ and every $\ell\ge1$,
\[
\trel\bigl(B(F,\ell)\bigr)\ge|E(F)|\,\ell+\frac{\ell^2}6+\frac13 .
\]
\end{lemma}

\begin{proof}
By \eqref{eq:dirichlet} and the self-adjointness of $P$ on $\ell^2(\pi)$, every nonzero $f$ with $\sum_vd_vf_v=0$ satisfies
\[
1-\lambda_2\le\frac{\sum_{uv\in E}(f_u-f_v)^2}{\sum_vd_vf_v^2}.
\]
Let $f=-1$ on the first copy of $F$, $f=1$ on the second, and $f=-1+2j/\ell$ at the $j$-th internal vertex of the path, counted from the first copy. The automorphism exchanging the two copies and reversing the path preserves degrees and changes $f$ into $-f$, so $\sum_vd_vf_v=0$. Only the $\ell$ path edges contribute to the numerator, each with $(2/\ell)^2$, so the numerator is $4/\ell$. In the denominator each copy of $F$ contributes its degree sum $2|E(F)|+1$, the copy of $o$ having gained one edge, and the internal vertices have degree $2$ and contribute
\[
2\sum_{j=1}^{\ell-1}\Bigl(1-\frac{2j}\ell\Bigr)^2=\frac{2(\ell-1)(\ell-2)}{3\ell}.
\]
Therefore
\[
\trel\ge\frac\ell4\Bigl(4|E(F)|+2+\frac{2(\ell-1)(\ell-2)}{3\ell}\Bigr)=|E(F)|\,\ell+\frac{\ell^2}6+\frac13 .\qedhere
\]
\end{proof}

\begin{proof}[Proof of Theorem~\ref{thm:main}]
The upper bound is Theorem~\ref{thm:capacity} with $\gamma=1$, together with $\trel\le K$. For the lower bound let $n\ge6$, $a=\lfloor n/3\rfloor\ge2$, $\ell=n-2a+1$, and $B_n=B(K_a,\ell)$, which has $n$ vertices. Since $(n-2)/3\le a\le n/3$, we have $a-1\ge(n-5)/3>0$ and $\ell\ge n/3+1$. Lemma~\ref{lem:barbell} gives
\[
\trel(B_n)\ge\frac{a(a-1)\ell}2+\frac{\ell^2}6+\frac13\ge\frac{(n-2)(n-5)(n+3)}{54}+\frac{(n+3)^2}{54}+\frac13=\frac{n^3-3n^2-5n+57}{54},
\]
which implies \eqref{eq:lower}. Since $K\ge\trel$, the upper and lower bounds together give $\max_GK(G)\sim\max_G\trel(G)\sim n^3/54$.
\end{proof}

\begin{proof}[Proof of Theorem~\ref{thm:clique}]
The upper bound is Theorem~\ref{thm:capacity} with $\gamma=1-1/r$. For the lower bound let $n\ge6r$, $m=\lfloor n/(3r)\rfloor\ge2$, $a=rm$, $\ell=n-2a+1$, and let $F$ be the complete $r$-partite graph with $r$ parts of size $m$, which is connected and has $(1-1/r)a^2/2$ edges. The barbell $B(F,\ell)$ has $n$ vertices. Its chromatic number is at most $r$: colour the first copy of $F$ properly with $r$ colours, colour the path alternately with the colour of its first vertex and one other colour, and colour the second copy of $F$ properly so that its copy of $o$ receives the colour already given to the last vertex of the path, which is possible after permuting the colours. Since $a\ge n/3-r$ and $\ell\ge n/3+1$, Lemma~\ref{lem:barbell} gives
\[
\trel\bigl(B(F,\ell)\bigr)\ge\Bigl(1-\frac1r\Bigr)\frac{a^2\ell}2\ge\Bigl(1-\frac1r\Bigr)\frac12\Bigl(\frac{n^2}9-\frac{2rn}3\Bigr)\frac n3=\Bigl(1-\frac1r\Bigr)\frac{n^3}{54}-\frac{(r-1)n^2}9,
\]
where we used $(n/3-r)^2\ge n^2/9-2rn/3\ge0$ for $n\ge6r$. A graph of chromatic number at most $r$ is $K_{r+1}$-free, so the upper bound applies to both classes, and the statements about the maxima follow. For $r=2$ the $K_3$-free graphs are the triangle-free graphs, and the graphs of chromatic number at most $2$ are the bipartite graphs.
\end{proof}

\begin{remark}\label{rem:scope}
Theorem~\ref{thm:main} determines the leading constant, not the extremal graphs for each $n$. The error term $O(n^{5/2})$ comes from the threshold $D=\sqrt n$ and has not been optimized.
\end{remark}

\section*{Acknowledgements and funding}
The research presented in this paper was supported by the European Research
Council (ERC) under the European Union's Horizon Europe research and
innovation programme (grant agreement No.~101041711), by the Simons
Foundation as part of the Collaboration on the Mathematical and Scientific
Foundations of Deep Learning, by Heights Labs, by Convex Nexus Capital, by
the Israel Science Foundation (grant number 2258/19), and by the Israel
Science Foundation (ISF Grant 4101/25). Computational resources were
provided in part by Google Cloud through its education credits program
(credit 530550103).

\section*{Declaration of competing interest}
The author holds equity in Heights Labs and in Convex Nexus Capital and
has a management role at Convex Nexus Capital. These entities are included
in the funding declaration above.

\section*{Declaration of generative AI and AI-assisted technologies}
The argument of Section~\ref{sec:upper} in the case $\gamma=1$ and the lower bound of Theorem~\ref{thm:main} were produced by OpenAI's GPT-6.1 model, run through the Codex agent. The capacity formulation and the multipartite barbells were produced earlier with OpenAI's GPT-5.6. Anthropic's Claude Opus~5.5 checked the arguments independently, re-derived every step, extended the upper-bound argument to the capacity condition~\eqref{eq:capacity}, and wrote the manuscript. These checks are not peer review. The author directed the work and is responsible for the content of the paper and for decisions concerning its dissemination.

\section*{Data availability}
No data were used.

\begin{thebibliography}{9}
\bibitem{AF}
D.~Aldous and J.~A.~Fill, \emph{Reversible Markov Chains and Random Walks on Graphs}, unfinished monograph, 2002; recompiled 2014. \url{https://www.stat.berkeley.edu/~aldous/RWG/book.html}.
\bibitem{ACTT}
S.~G.~Aksoy, F.~Chung, M.~Tait and J.~Tobin, \emph{The maximum relaxation time of a random walk}, Adv. in Appl. Math. \textbf{101} (2018), 1--14. \href{https://doi.org/10.1016/j.aam.2018.07.002}{doi:10.1016/j.aam.2018.07.002}.
\bibitem{Breen}
J.~Breen, S.~Butler, N.~Day, C.~DeArmond, K.~Lorenzen, H.~Qian and J.~Riesen, \emph{Computing Kemeny's constant for a barbell graph}, Electron. J. Linear Algebra \textbf{35} (2019), 583--598. \href{https://doi.org/10.13001/ela.2019.5175}{doi:10.13001/ela.2019.5175}.
\bibitem{Faught}
N.~Faught, M.~Kempton and A.~Knudson, \emph{A 1-separation formula for the graph Kemeny constant and Braess edges}, J. Math. Chem. \textbf{60} (2022), 49--69. \href{https://doi.org/10.1007/s10910-021-01294-8}{doi:10.1007/s10910-021-01294-8}.
\bibitem{MS}
T.~S.~Motzkin and E.~G.~Straus, \emph{Maxima for graphs and a new proof of a theorem of Tur\'an}, Canad. J. Math. \textbf{17} (1965), 533--540. \href{https://doi.org/10.4153/CJM-1965-053-6}{doi:10.4153/CJM-1965-053-6}.
\end{thebibliography}
\end{document}
